% アクター指向言語における「1アクター = 1実行担体」モデルの対応表
%   lualatex -interaction=nonstopmode actor_process_models_ja.tex   （2回）
\documentclass[11pt,a4paper]{ltjsarticle}
\usepackage{amsmath,amssymb}
\usepackage{listings}
\usepackage{xcolor}
\usepackage{geometry}
\usepackage{array}
\usepackage{booktabs}
\usepackage{longtable}
\usepackage{pdflscape}
\usepackage[hidelinks]{hyperref}
\geometry{margin=20mm}

\definecolor{codebg}{gray}{0.96}
\lstset{
  basicstyle=\ttfamily\small,
  backgroundcolor=\color{codebg},
  frame=single,
  framesep=4pt,
  rulecolor=\color{gray},
  breaklines=true,
  showstringspaces=false,
  columns=fullflexible,
  keepspaces=true,
  commentstyle=\itshape\color{gray!130},
}
\lstnewenvironment{term}{\lstset{language={},basicstyle=\ttfamily\footnotesize}}{}

\newcommand{\dg}{\textsuperscript{\dag}}
\newcommand{\yes}{\textbf{○}}
\newcommand{\no}{×}

\title{アクター指向言語における\\「1 アクター $=$ 1 実行担体」モデルの対応表\\
{\large ---\ Erlang の「プロセス」・OS スレッド・軽量プロセスの区別と、
AIPL/ABCL$^{c}$+ の位置づけ ---}}
\author{ABCL$^{c}$+ / AIOS 処理系（\texttt{yaskodama/abclcp}）}
\date{2026年07月31日}

\begin{document}
\maketitle

\begin{abstract}
「Erlang などのアクター言語に、アクター 1 つにつきプロセス 1 つに限定した
モデルはあるか」という問いに答える。結論を先に書くと、
\textbf{ある。ただし「プロセス」の語が指すものを 3 層に分けないと問いが定まらない。}
Erlang は確かに「1 アクター $=$ 1 プロセス」だが、その\emph{プロセス}は
OS のプロセスでもスレッドでもなく BEAM の軽量プロセスである。
一方、\textbf{1 アクターを 1 本の OS スレッドに固定する}処理系も、
\textbf{1 アクターを 1 個の OS プロセスに固定する}処理系も実在する。
前者の代表は Celluloid・Pykka・AmbientTalk・本処理系 AIPL、
後者の代表は Ray と Thespian である。
Akka と CAF は既定では多重化しつつ、\emph{オプトインで}1 対 1 に切り替えられる。

本書は主要なアクター指向言語・処理系を横断して、
(1) アクターが載る実行担体、(2) 1 対 1 に専有するか、
(3) アクター本体がブロックできるか、(4) 実用上のアクター数、
の 4 点で対応表を作る（\S\ref{sec:table}）。
そのうえで、なぜ 1 対 1 の OS スレッド束縛が稀なのか（コスト、\S\ref{sec:cost}）、
それでもなぜ選ぶのか（同期待ちが書けること、\S\ref{sec:why}）を述べ、
最後に AIPL の設計上の位置づけを示す（\S\ref{sec:aipl}）。

AIPL は 1 アクターにつき 1 本の OS スレッドを生成する。
本書のために実測したところ、
\textbf{4000 アクターまでは生成でき、5000 で \texttt{Thread.create:
Resource temporarily unavailable} で失敗した}。
上限は macOS の \texttt{kern.num\_taskthreads} $=4096$ である。
この数字がこのモデルの代償を端的に表している。
\end{abstract}

\tableofcontents

\section{問いの整理 --- 「プロセス」は 3 層ある}

アクター言語の文献で「プロセス」と呼ばれるものには、
コストが 3 桁以上違う 3 つの層がある。この区別をしないと
「1 アクター 1 プロセス」という言明は意味を持たない。

\begin{center}
\small
\begin{tabular}{@{}p{3.4cm}p{5.0cm}p{4.4cm}l@{}}
\toprule
層 & 何か & 1 個あたりの費用 & 実用上の個数 \\
\midrule
OS プロセス
  & 独立したアドレス空間。カーネルが管理
  & 数 MB $\sim$、生成 ms オーダー & $10^{2}$ \\
OS スレッド (pthread)
  & アドレス空間を共有。カーネルが管理・プリエンプト
  & スタック 512\,KB (macOS) $\sim$ 8\,MB (Linux) & $10^{3}$ \\
軽量プロセス／ グリーンスレッド
  & 言語ランタイムが管理する私有スタック。M:N で OS スレッドへ載せる
  & BEAM で 326 語（$\approx$2.6\,KB）& $10^{6}$ \\
ターン（スタックを 持たない）
  & 「1 メッセージ $=$ 1 関数呼び出し」。実行の間だけプールのスレッドを借りる
  & Pony のアクターで $\approx$240\,B & $10^{6\sim7}$ \\
\bottomrule
\end{tabular}
\end{center}

\textbf{Erlang の「プロセス」は 3 段目である。}
Erlang は確かに「1 アクター $=$ 1 プロセス」を厳密に守るが、
それは OS から見れば何でもない。BEAM はコア数と同じだけの
スケジューラスレッドを立て、その上に $10^{6}$ 個の軽量プロセスを
M:N で載せる。したがって
「Erlang はアクターごとにプロセスを持つから重いのでは」という直感は誤りで、
むしろ\emph{軽いからこそ}1 対 1 を貫ける。

問いを言い直すと、次の 2 つになる。

\begin{enumerate}
\item \textbf{OS の実行資源（スレッド／プロセス）を 1 アクターに専有させる
  処理系はあるか。} $\rightarrow$ ある（\S\ref{sec:oneone}）。
\item \textbf{それを言語の既定にしている処理系はあるか。} $\rightarrow$
  ごく少数だが、ある。AIPL はその 1 つである。
\end{enumerate}

\section{対応表の読み方}

対応表では次の 4 軸で分類する。

\begin{description}
\item[実行担体] アクターが実行されるとき、何の上で走るか。
  OS プロセス／OS スレッド／軽量プロセス／プールのターン。
\item[1 対 1 か] その担体をアクターが\emph{専有}するか、
  多数のアクターで多重化するか。\yes\ は専有、\no\ は多重化。
\item[ブロック可] アクターのメソッド本体の\emph{途中}で
  同期的に待てるか。待つには私有スタックが要る。
  待てない処理系（run-to-completion）では、
  同期呼び出しは継続・コールバック・\texttt{await} に書き換える必要がある。
\item[実用規模] 1 ノードで無理なく作れるアクター数の桁。
\end{description}

この 4 軸を選んだのは、\textbf{第 3 軸が他の 3 つを決めてしまう}からである。
アクター本体の途中で待てるようにすると、待っている間の呼び出し状態を
保持する場所が要る。その場所として OS スレッドのスタックを使えば
第 1 軸と第 2 軸が「OS スレッド・専有」に決まり、
その結果として第 4 軸がカーネルの上限で頭打ちになる。
逆に待てない設計（run-to-completion）にすれば担体を専有する必要が無くなり、
アクターは単なるオブジェクトになって $10^{6}$ 個作れる。
つまりこの表は、実質的には\textbf{「同期的に待てるか」という 1 つの設計判断が
どこまで波及するか}の一覧である。この点は \S\ref{sec:why} で改めて述べる。

なお「メールボックスの選択受信ができるか」「監視ツリーがあるか」
「メッセージ順序をどこまで保証するか」といった軸もありうるが、
これらは担体の選択とは概ね独立に決められるので、本表では扱わない
（AIPL と Erlang の比較としては \S\ref{sec:aipl} の末尾に挙げた）。

\begin{quote}\small
凡例：\dg\ を付けた行は、標準的な文献・処理系の知識にもとづく記述で、
本書の作成にあたって一次資料での再確認を\emph{していない}ものである。
\dg\ の無い行は \S\ref{sec:src} に挙げた資料で確認した。
AIPL の行は実装のソースと実測による。
\end{quote}

\begin{landscape}
\section{対応表}
\label{sec:table}

\subsection{アクターが言語の構文である言語}

\begin{longtable}{@{}p{3.3cm}p{4.6cm}>{\raggedright\arraybackslash}p{2.5cm}>{\raggedright\arraybackslash}p{2.6cm}>{\raggedright\arraybackslash}p{1.5cm}p{6.0cm}@{}}
\toprule
言語（初出） & 実行担体 & 1 対 1 か & ブロック可 & 実用規模 & 備考 \\
\midrule
\endfirsthead
\toprule
言語（初出） & 実行担体 & 1 対 1 か & ブロック可 & 実用規模 & 備考 \\
\midrule
\endhead
\bottomrule
\endfoot

Act1 / Act2 / Act3 (1981--)\dg
  & 概念上のみ。実装は Lisp 上のスケジューラ
  & ---
  & ---
  & ---
  & 理論的なアクターモデルの原型。担体の議論をしていない \\

ABCL/1 (1986)\dg
  & 「各オブジェクトが固有の活動を持つ」。実装は Lisp の軽量プロセス
  & 概念上 \yes
  & \yes\ (\texttt{now})
  & $10^{3}$
  & \texttt{past}/\texttt{now}/\texttt{future} の 3 送信。AIPL の直接の祖先 \\

Concurrent Smalltalk (1987)\dg
  & 処理系の軽量プロセス
  & \no
  & \yes
  & $10^{4}$
  & \\

\textbf{Erlang} (1986/98)
  & \textbf{BEAM 軽量プロセス}。コア数個のスケジューラスレッドへ M:N
  & \yes\ （軽量プロセスと）
  & \yes\ (\texttt{receive})
  & $10^{6}$
  & 新規プロセスは 326 語。既定の上限 \texttt{+P 262144}。選択的受信あり。
    \textbf{問いの「1 アクター 1 プロセス」はここでは成立するが、
    プロセス $=$ 軽量プロセス} \\

Elixir / Gleam / LFE\dg
  & 同上（BEAM を共有）
  & \yes
  & \yes
  & $10^{6}$
  & 実行モデルは Erlang と同一 \\

E (1997)
  & \textbf{vat} $=$ 1 OS スレッド。1 つの vat に多数のオブジェクトが載る
  & \no
  & \no
  & ---
  & 1 vat に $N$ 個のオブジェクトが載るので、担体はアクター単位ではない。
    非同期送信のみで、同期呼び出しと構文で区別する \\

AmbientTalk (2006)
  & アクター $=$ イベントループ $=$ 1 JVM スレッド
  & \yes
  & \no
  & $10^{2\sim3}$
  & E の vat モデルを継承。
    \textbf{1 アクター 1 OS スレッドの言語レベルの例} \\

SALSA (1999)
  & 初期は 1 アクター 1 Java スレッド。SALSA Lite で
    ステージ（スレッド）に多数のアクターを載せる方式へ
  & 初期 \yes\ $\rightarrow$ \no
  & \yes\dg
  & $10^{3}\rightarrow10^{5}$
  & \textbf{同一言語がスレッド専有から多重化へ移った実例} \\

Scala Actors (2006, 2.11 で非推奨)
  & \texttt{receive} $=$ スレッドを占有 /
    \texttt{react} $=$ 継続クロージャのみ
  & \textbf{選べる}
  & \texttt{receive} なら \yes
  & \texttt{react} で $10^{5}$
  & \textbf{同じ言語の中に両モデルが並置された唯一に近い例}。
    Haller \& Odersky「Actors that Unify Threads and Events」 \\

Io (2002)\dg
  & コルーチン（言語ランタイムが管理）
  & \yes\ （コルーチンと）
  & \yes
  & $10^{4}$
  & オブジェクトが動的にアクター化する \\

Pony (2015)
  & 担体を持たない。behaviour は run-to-completion で
    コア数個のスケジューラが work-stealing
  & \no
  & \no
  & $10^{6\sim7}$
  & ブロックを設計として禁じている。空アクター $\approx$240\,B。
    参照能力型でデータ競合が起こらないことを保証 \\

Encore (2015)\dg
  & Pony ランタイム上の active object
  & \no
  & \no\ (future)
  & $10^{6}$
  & 能力型 $+$ 並列コンビネータ \\

Swift \texttt{actor} (2021)
  & Executor 上の Job。既定は協調スレッドプール（コア数程度）
  & \no
  & \no\ (\texttt{await})
  & $10^{6}$
  & カスタム executor でスレッド固定は\emph{可能}だが既定ではない。
    再入（reentrancy）を許す \\

Motoko (2019)
  & canister $=$ アクター。1 メッセージずつ処理
  & \yes\ (canister)
  & \no\ (\texttt{await})
  & ---
  & \texttt{await} がメッセージを分割し commit point になる。
    待ちが状態のコミット境界になるという珍しい意味論 \\

Rebeca (2001)\dg
  & モデル検査用。各 rebec が固有の制御を持つ
  & 概念上 \yes
  & ---
  & ---
  & 実行より検証が目的。Timed Rebeca は時間も持つ \\

D \texttt{std.concurrency} (2010)\dg
  & \texttt{spawn} が \textbf{OS スレッド}を作る
  & \yes
  & \yes\ (\texttt{receive})
  & $10^{3}$
  & 標準ライブラリだが言語と密着。
    \textbf{1 アクター 1 OS スレッドの例} \\

\textbf{AIPL / ABCL$^{c}$+} (2025--26)
  & \textbf{OS スレッド (pthread)}。\texttt{new} ごとに \texttt{Thread.create}
  & \textbf{\yes（既定かつ唯一）}
  & \textbf{\yes\ (\texttt{now})}
  & \textbf{実測 4096}
  & OCaml 5.1 の \texttt{Thread}。同一ドメイン内なので
    \emph{並行だが並列ではない}。\S\ref{sec:aipl} \\
\end{longtable}

\subsection{ライブラリ・フレームワーク}

\begin{longtable}{@{}p{3.3cm}p{4.6cm}>{\raggedright\arraybackslash}p{2.5cm}>{\raggedright\arraybackslash}p{2.6cm}>{\raggedright\arraybackslash}p{1.5cm}p{6.0cm}@{}}
\toprule
処理系（言語） & 実行担体 & 1 対 1 か & ブロック可 & 実用規模 & 備考 \\
\midrule
\endfirsthead
\toprule
処理系（言語） & 実行担体 & 1 対 1 か & ブロック可 & 実用規模 & 備考 \\
\midrule
\endhead
\bottomrule
\endfoot

Akka / Akka.NET (JVM/.NET)
  & 既定は共有スレッドプール（dispatcher、8--64 スレッド）。
    \textbf{\texttt{PinnedDispatcher} でアクター専用スレッド}
  & 既定 \no ／可 \yes
  & \no\ (\texttt{ask})
  & $10^{6}$
  & run-to-completion なので待てない。同期呼び出しは \texttt{ask} が返す
    Future で書く。「重要なアクターに専用スレッドを与える」
    「バルクヘッドを作る」ための機能として 1 対 1 が用意されている \\

CAF (C++)
  & 既定は協調スケジューリング（数百バイト）。
    \textbf{\texttt{spawn<detached>} で専有スレッド}
  & 既定 \no ／可 \yes
  & detached なら \yes
  & $10^{6}$
  & ブロッキング I/O を呼ぶアクターを他から隔離する用途。
    \textbf{「ブロックしたいなら専有スレッド」という対応関係が明示的} \\

Orleans (.NET)\dg
  & 仮想アクター（grain）。ターン単位でスレッドプール上に活性化
  & \no
  & \no\ (\texttt{await})
  & $10^{6}$
  & 使われていない grain は非活性化される。担体を持たないので当然 \\

Proto.Actor / Dapr Actors\dg
  & 同上（ターンベース）
  & \no
  & \no
  & $10^{6}$
  & \\

Actix (Rust)\dg
  & Arbiter（$=$1 スレッド）に多数のアクター。
    専用 Arbiter を与えれば 1 対 1
  & 既定 \no ／可 \yes
  & \no
  & $10^{6}$
  & \\

Ractor / Bastion (Rust)\dg
  & tokio タスク／ランタイムの軽量プロセス
  & \no
  & \no
  & $10^{6}$
  & Bastion は Erlang の監視ツリーを模す \\

Cloud Haskell\dg
  & \texttt{forkIO} 軽量スレッド（HEC へ M:N）
  & \yes\ （軽量スレッドと）
  & \yes
  & $10^{6}$
  & Erlang に最も近い構図 \\

\textbf{Ray} (Python)
  & \textbf{1 アクター $=$ 1 専有 OS プロセス（ワーカー）}
  & \textbf{\yes（OS プロセス）}
  & \yes
  & $10^{2\sim3}$／ノード
  & タスクと違いワーカーを再利用しない。
    \textbf{問いに対する最も直接的な「はい」} \\

\textbf{Thespian} (Python)
  & \texttt{multiprocTCPBase} 等では \textbf{1 アクター 1 OS プロセス}。
    \texttt{simpleSystemBase} は単一プロセス
  & \textbf{選べる}
  & \yes
  & $10^{2}$
  & 同一のアクターコードのまま担体だけ差し替えられる設計 \\

\textbf{Pykka} (Python)
  & \textbf{\texttt{ThreadingActor} $=$ 1 \texttt{threading.Thread}}
  & \textbf{\yes}
  & \yes
  & $10^{3}$
  & 「各 ThreadingActor は通常のスレッドで実行される」 \\

\textbf{Celluloid} (Ruby)
  & \textbf{1 アクター $=$ 1 Ruby スレッド}。
    メソッド呼び出しごとに Fiber で待ちを退避
  & \textbf{\yes}
  & \yes
  & $10^{3}$
  & 「2000 個以上のアクターが作れない」という報告が実際にある。
    \textbf{このモデルの代償がそのまま現れた例} \\

Nact / Comedy (JS)\dg
  & 単一のイベントループ
  & \no
  & \no
  & $10^{5}$
  & \\
\end{longtable}

\subsection{参考：アクターではないが対比される仕組み}

\begin{longtable}{@{}p{3.3cm}p{4.6cm}>{\raggedright\arraybackslash}p{2.5cm}>{\raggedright\arraybackslash}p{2.6cm}>{\raggedright\arraybackslash}p{1.5cm}p{6.0cm}@{}}
\toprule
 & 実行担体 & 1 対 1 か & ブロック可 & 実用規模 & 備考 \\
\midrule
\endfirsthead
\endhead
\bottomrule
\endfoot

Go (goroutine $+$ channel)\dg
  & goroutine（初期スタック 2\,KB、M:N）
  & 慣用的に \yes
  & \yes
  & $10^{6}$
  & CSP であってアクターではない（宛先が名前でなくチャネル）。
    「1 状態 $=$ 1 goroutine」の書き方が事実上アクター \\

Vert.x (JVM)\dg
  & verticle をイベントループスレッドに固定。1 loop : $N$ verticle
  & \no
  & \no
  & $10^{5}$
  & \\

Java 21 virtual threads\dg
  & JVM 管理の軽量スレッド
  & ---
  & \yes
  & $10^{6}$
  & \textbf{「1 アクター 1 スレッド」の経済性を変える}。
    \S\ref{sec:why} \\
\end{longtable}
\end{landscape}

\section{「1 アクター $=$ 1 OS 実行資源」を採る処理系（抽出）}
\label{sec:oneone}

対応表から該当するものだけを抜き出すと、次のようになる。

\begin{center}
\small
\begin{tabular}{@{}llp{7.4cm}@{}}
\toprule
 & 担体 & 処理系 \\
\midrule
\textbf{既定でそう} & OS プロセス & Ray（Python）、Thespian（multiproc 系） \\
                    & OS スレッド & \textbf{AIPL/ABCL$^{c}$+}、AmbientTalk、
                                    Celluloid（Ruby）、Pykka \texttt{ThreadingActor}、
                                    D \texttt{std.concurrency}\dg、
                                    初期の SALSA \\
\midrule
\textbf{オプトインで} & OS スレッド & Akka \texttt{PinnedDispatcher}、
                                      CAF \texttt{spawn<detached>}、
                                      Actix の専用 Arbiter\dg、
                                      旧 Scala Actors の \texttt{receive}、
                                      Swift のカスタム executor \\
\midrule
\textbf{軽量プロセスと 1 対 1} & 言語ランタイム
                    & Erlang / Elixir / Gleam / LFE、Cloud Haskell\dg、Io\dg \\
\midrule
\textbf{担体を専有しない} & プールのターン
                    & Pony、Encore、Swift（既定）、Akka（既定）、CAF（既定）、
                      Orleans、Proto.Actor、Dapr、Nact \\
\bottomrule
\end{tabular}
\end{center}

読み取れることは 3 つある。

\begin{enumerate}
\item \textbf{OS スレッド専有は、言語の既定としては少数派である。}
  該当するのは研究用言語（AmbientTalk、初期 SALSA、AIPL）と、
  動的言語のライブラリ（Celluloid、Pykka）に偏る。
\item \textbf{大規模実用系はすべてオプトインにしている。}
  Akka と CAF は、既定では多重化しつつ
  「ブロッキング I/O を呼ぶアクターだけ隔離する」ために
  専有スレッドを用意している。つまり\emph{専有は例外的な道具}という位置づけである。
\item \textbf{分散フレームワークでは 1 アクター 1 OS プロセスが主流になる。}
  Ray と Thespian がそうで、これは並列性（GIL の回避）と
  障害隔離のためであって、アクターモデルの要請ではない。
\end{enumerate}

\section{なぜ稀か --- コスト}
\label{sec:cost}

\begin{center}
\small
\begin{tabular}{@{}lp{6.2cm}l@{}}
\toprule
担体 & 1 個あたり & 出典・備考 \\
\midrule
BEAM 軽量プロセス & 326 語（初期ヒープ 233 語を含む）
  & Erlang 効率化ガイド \\
Pony のアクター & $\approx$240\,B（64bit の空アクター）
  & Pony FAQ \\
goroutine & 初期スタック 2\,KB\dg & \\
OS スレッド & スタック 512\,KB (macOS) $\sim$ 8\,MB (Linux)、
              加えてカーネル側の資源 & \\
OS プロセス & 数 MB $\sim$ & \\
\bottomrule
\end{tabular}
\end{center}

差は容量よりも\textbf{カーネルが管理する資源には上限がある}ことに現れる。
本書のために AIPL で実測した結果を挙げる。

\begin{center}
\small
\begin{tabular}{@{}lll@{}}
\toprule
アクター数 & 生成の所要時間（プロセス起動込み） & 結果 \\
\midrule
1\,000  & 0.32 秒 & 成功 \\
2\,000  & 1.03 秒 & 成功 \\
3\,000  & 2.14 秒 & 成功 \\
4\,000  & 3.66 秒 & 成功 \\
5\,000  & --- & \textbf{\texttt{Thread.create: Resource temporarily unavailable}} \\
\bottomrule
\end{tabular}
\end{center}

\begin{term}
$ sysctl kern.num_taskthreads
kern.num_taskthreads: 4096
\end{term}

上限は明らかに OS のプロセスあたりスレッド数上限である。
すなわち\textbf{このモデルではアクター数の上限が言語ではなくカーネルで決まる}。
所要時間も超線形に伸びている（1000 個で 0.32 秒、4000 個で 3.66 秒）。

同じ現象は他の処理系でも観測されていて、Celluloid には
「自分のマシンでは 2000 個以上のアクター／スレッドを作れない」という
報告がそのまま issue として残っている。
Erlang が既定で $2^{18}=262\,144$ プロセスを許し、
実測で $10^{6}$ を扱えるのとは 2 桁以上違う。

\section{ではなぜ選ぶのか --- ブロックできること}
\label{sec:why}

コストが 2 桁以上不利なのに 1 アクター 1 スレッドを選ぶ理由は 1 つである。
\textbf{アクター本体の途中で同期的に待てるようにするため}である。

メソッドの途中で待つには、待っている間の\emph{呼び出し状態}をどこかに
置かねばならない。置き場所は 3 つしかない。

\begin{enumerate}
\item \textbf{OS スレッドのスタック} --- 何も書き換えずに待てる。
  代わりにアクター数が $10^{3}$ で頭打ちになる。
  AIPL、Celluloid、Pykka、CAF の detached、Akka の PinnedDispatcher。
\item \textbf{言語ランタイムの軽量プロセスのスタック} --- 待てて、かつ安い。
  ただしランタイムを自前で持つ必要がある。Erlang、Go、Java 21 の virtual thread。
\item \textbf{継続（クロージャ）} --- スタックを持たない。
  代わりに\emph{待つ構文を書けない}ので、
  \texttt{await} やコールバックへの書き換えが必要になる。
  Pony、Swift、Akka、Orleans、Motoko。
\end{enumerate}

この三択が同じ言語の中に並置された例が、旧 Scala Actors の
\texttt{receive}（スレッドを占有して待つ）と
\texttt{react}（継続だけ残して待つ）である。
Haller と Odersky の「Actors that Unify Threads and Events」は
まさにこの選択を統一しようとした仕事で、
「アクターはフルのスレッドスタックで中断することも、
継続クロージャだけで中断することもできる」と述べている。
第 1 の形がスレッドベース、第 2 の形がイベントベースである。

CAF の \texttt{detached} も同じ対応関係を明示している。
「ブロッキングなシステムコールを使うアクターはスレッドを止めて
不均衡や飢餓を招くので、協調スケジューリングから外す」という説明で、
\textbf{ブロックしたいなら専有スレッド}という交換条件がそのまま書かれている。

\begin{quote}\small
\textbf{Java 21 の virtual thread は、この経済性を変えつつある。}
選択肢 2 が JVM の標準機能として使えるようになったので、
「待てる」ことと「安い」ことが両立する。
同じことは Go では以前から成り立っていた。
OS スレッド専有が正当化される場面は、
今後は「ランタイムを持たない／持てない言語」に限られていく可能性が高い。
\end{quote}

\section{AIPL / ABCL$^{c}$+ の位置づけ}
\label{sec:aipl}

\subsection{何をしているか}

AIPL は \texttt{new C(...)} のたびに \texttt{Thread.create} を呼び、
アクターごとに 1 本の OCaml スレッドを立てる
（\texttt{src/eval\_thread.ml}）。
アクターループはメールボックスからメッセージを取り出して逐次に処理し、
\texttt{now} での待ちは条件変数でブロックする。

\begin{lstlisting}[language={}]
ignore (Thread.create (fun () -> actor_loop actor_inst) ());
\end{lstlisting}

正確には次の 2 点に注意が要る。

\begin{itemize}
\item OCaml 5.1 の \texttt{Thread} は\textbf{本物の OS スレッド}である。
  したがってアクター数の上限は OS のスレッド数上限になる（実測 4096）。
\item ただし本処理系は \texttt{Domain} を使っていないので、
  すべてのスレッドが単一ドメインに属する。
  OCaml のランタイムロックにより、\textbf{同時に OCaml コードを実行できるのは
  1 本だけ}である。すなわち AIPL のアクターは
  \emph{並行 (concurrent) だが並列 (parallel) ではない}。
  コストは OS スレッドのぶん払っているが、
  マルチコアの利得は現状では得ていない。
\end{itemize}

\subsection{なぜこの設計なのか}

AIPL は ABCL/1 の \texttt{now}（同期送信）を言語の中心に据えている。

\begin{lstlisting}[language={}]
var left = now stock.take(k) timeout 100 else 0;
if (left > 0) { reply("ok, left=" ++ left); } else { reply("sold out"); }
\end{lstlisting}

これは\textbf{メソッド本体の途中で待つ}構文である。
\S\ref{sec:why} の三択でいえば選択肢 1 を採ったことになり、
だから 1 アクター 1 OS スレッドになる。
Pony や Swift のように継続へ分割する方式を採れば安くなるが、
そのときは \texttt{now} という構文自体が書けなくなる
（\texttt{await} と継続に置き換わる）。

\subsection{代償が言語仕様に現れている}

このモデルの代償は、AIPL では 2 か所に見える形で現れている。

\begin{enumerate}
\item \textbf{期限付きの待ちを導入せざるを得なかった。}
  アクターが 1 台 1 スレッドで、待ちがそのスレッドをブロックする以上、
  返信が来なければそのアクターは永久に停止する。
  そこで \texttt{now ... timeout $n$ else $e$} を入れ、
  待ちに上限を置いた。期限付きの待ちだけを使うプログラムでは
  デッドロック自由が定理として言える。
  \textbf{これは実行モデルの制約が型システムの設計を決めた例である。}
\item \textbf{\texttt{now self.m()} を文法で禁止した。}
  自分自身への同期送信は、自分がそのメソッドを処理中である以上
  返信を作る主体がいないので必ず自己デッドロックになる。
  1 アクター 1 スレッドだからこそ確実にそうなるので、構文として書けなくした。
\end{enumerate}

\subsection{他との比較でどこに立つか}

\begin{center}
\small
\begin{tabular}{@{}lll@{}}
\toprule
 & AIPL & 対照 \\
\midrule
担体          & OS スレッド専有 & Erlang: 軽量プロセス専有 \\
待てるか      & \yes（\texttt{now}）& Pony/Swift: \no \\
上限          & 4096（カーネル）& Erlang: $10^{6}$ \\
並列性        & \no（単一ドメイン）& Erlang/Pony: \yes（コア数） \\
選択的受信    & \yes（\texttt{select}）& Akka: \no（stash で代用）\\
待ちの期限    & \yes（言語機能）& 多くは API レベル（\texttt{ask} のタイムアウト）\\
\bottomrule
\end{tabular}
\end{center}

対応表の中で AIPL に最も近いのは
\textbf{AmbientTalk（1 アクター 1 JVM スレッド）}と
\textbf{Celluloid（1 アクター 1 Ruby スレッド）}だが、
両者と違って AIPL は\emph{同期待ちを言語の一級の構文として持ち、
かつその期限を型検査で見る}。
「1 アクター 1 OS スレッド」＋「同期待ちが書ける」＋
「待ちの期限が静的に検査される」という組み合わせは、
調べた範囲では他に見当たらない。

\section{まとめ}

\begin{enumerate}
\item 「1 アクター $=$ 1 プロセス」は Erlang で成立するが、
  その\emph{プロセス}は BEAM の軽量プロセスであって OS の資源ではない。
\item OS の実行資源を 1 アクターに専有させるモデルは実在する。
  OS プロセス単位では Ray と Thespian、
  OS スレッド単位では AmbientTalk・Celluloid・Pykka・初期 SALSA・
  D \texttt{std.concurrency}・AIPL。
  Akka と CAF はオプトインとして持っている。
\item 言語の\emph{既定}としてそうしている例は少数で、
  研究用言語と動的言語のライブラリに偏る。
  理由は単純にコストで、アクター数の上限がカーネルで決まってしまう
  （AIPL の実測で 4096）。
\item それでも選ぶ理由は「メソッド本体の途中で同期的に待てる」ことに尽きる。
  待ち状態の置き場所がスレッドスタック／軽量プロセス／継続の三択であり、
  ランタイムを持たない処理系は第 1 の選択肢しか採れない。
\item AIPL はこの第 1 の選択肢を採った。その代償が
  「待ちに期限を書かせる」という言語仕様として表に出ている。
\end{enumerate}

\section{出典}
\label{sec:src}

本文中で確認に用いた資料。

\begin{itemize}\small
\item Erlang System Documentation, \emph{Efficiency Guide --- Processes}（初期ヒープ 233 語、
  新規プロセス 326 語）\\
  \url{https://www.erlang.org/doc/system/eff_guide_processes.html}
\item Ray Documentation, \emph{Actors}（各アクターが専有 Python プロセスで動く）\\
  \url{https://docs.ray.io/en/latest/ray-core/actors.html}
\item Akka Documentation, \emph{Dispatchers}（\texttt{PinnedDispatcher} は
  アクターごとに 1 スレッドのプールを持つ）\\
  \url{https://doc.akka.io/libraries/akka-core/current/typed/dispatchers.html}
\item CAF Documentation, \emph{Actors}（協調スケジューリングと \texttt{detached}）\\
  \url{https://actor-framework.readthedocs.io/en/stable/core/Actors.html}
\item P. Haller, M. Odersky, \emph{Actors that Unify Threads and Events}, 2007
  （\texttt{receive} と \texttt{react}、スレッドスタックと継続クロージャ）\\
  \url{https://infoscience.epfl.ch/bitstreams/1b9cd67f-c907-4fac-8cfd-b0829b000dcf/download}
\item Swift Concurrency（Executor と協調スレッドプール、再入）\\
  \url{https://wwdcnotes.com/documentation/wwdc21-10254-swift-concurrency-behind-the-scenes/}
\item Internet Computer Documentation, \emph{Actors and async data}
  （canister $=$ アクター、\texttt{await} が commit point）\\
  \url{https://internetcomputer.org/docs/motoko/main/writing-motoko/actors-async}
\item AmbientTalk Tutorial, \emph{Concurrent Programming with Actors}
  （E の vat モデル、イベントループスレッド）\\
  \url{https://soft.vub.ac.be/amop/at/tutorial/actors}
\item Pony FAQ, \emph{Runtime}（スケジューラはコア数、空アクター $\approx$240--256\,B）\\
  \url{https://www.ponylang.io/faq/runtime/}
\item Pykka Documentation, \emph{Threading}（\texttt{ThreadingActor} は
  \texttt{threading.Thread} で実行される）\\
  \url{https://pykka.readthedocs.io/en/release-2.x/runtimes/threading/}
\item Thespian Documentation, \emph{Using Thespian}（system base による
  プロセス／スレッドの違い）\\
  \url{https://thespianpy.com/doc/using.html}
\item Celluloid（各アクターは専用スレッド、メソッド呼び出しは Fiber）\\
  \url{https://celluloid.io/} ／
  \url{https://github.com/celluloid/celluloid/issues/116}（2000 個で頭打ちの報告）
\item T. Desell, C. Varela, \emph{SALSA Lite: A Hash-Based Actor Runtime for
  Efficient Local Concurrency}, 2014（ステージ方式への移行）\\
  \url{https://link.springer.com/chapter/10.1007/978-3-662-44471-9_7}
\item AIPL: 本リポジトリ \texttt{src/eval\_thread.ml}（\texttt{Thread.create}）、
  および本書のための実測（macOS 26.5.2, OCaml 5.1.1,
  \texttt{kern.num\_taskthreads} $=4096$）
\end{itemize}

\end{document}
